\par\noindent\textbf{(1)} 井内波数 \(k=\sqrt{2m(E+V_0)}/\hbar\)，井外指数衰减常数 \(\kappa=\sqrt{-2mE}/\hbar\)。两者满足 \(k^2+\kappa^2=2mV_0/\hbar^2\)。对称势允许波函数取偶态 \(\cos kx\) 或奇态 \(\sin kx\)，外区分别按宇称接指数尾。
\par\noindent\textbf{(2)} 在 \(x=a\) 处波函数与导数连续。偶态内外对数导数匹配给 \(k\tan(ka)=\kappa\)，奇态给 \(-k\cot(ka)=\kappa\)。结合圆关系 \(k^2+\kappa^2=\mathrm{const}\) 可图解每个束缚能根；不能只用井内量子化 \(ka=n\pi\) 而忽略尾波。
\par\noindent\textbf{(3)} 最低偶态在 \(ka\in(0,\pi/2)\) 内的 \(k\tan(ka)\) 从零增大，而 \(\kappa=\sqrt{k_0^2-k^2}\) 从正变零，必有交点，因此任意小正 \(V_0\) 在一维都至少有一个束缚态。奇态最低根要求 \(ka>\pi/2\)，从阈值 \(\kappa\to0\) 得 \(k_0a>\pi/2\)。
\par\noindent\textbf{(4)} 浅阱 \(k_0a\ll1\) 下 \(ka\ll1\)，偶态方程给 \(\kappa\approx k^2a\approx(2mV_0/\hbar^2)a\)。所以 \(E\approx-\hbar^2\kappa^2/(2m)=-2mV_0^2a^2/\hbar^2\)。当 \(V_0\downarrow0\) 束缚能二次趋零，波函数外尾尺度 \(\kappa^{-1}\) 发散。
